\begin{knowbox}
\textbf{Nội dung:} Chương 1 -- 13 (tỉ số phần trăm, hình hộp chữ nhật và hình lập phương, thể tích, sơn bề mặt, chuyển động đều, thống kê và xác suất, sơ đồ Venn, tính ngược, giả thiết tạm). \quad \textbf{Thời gian gợi ý:} $60$ phút.
\end{knowbox}
\gap
\begin{probbox}{1}{\relax}
\begin{cols3}
\pa{a} Tìm $35\%$ của $240$. & \pa{b} $27$ bằng bao nhiêu phần trăm của $45$? & \pa{c} Biết $15\%$ của một số là $27$. Tìm số đó.
\end{cols3}
\end{probbox}
\probsol{
\pa{a} $240\times35:100=84$.\par
\pa{b} $27:45=0{,}6=60\%$.\par
\pa{c} $27:15\times100=180$.
\dapso{\pa{a} $84$;\ \pa{b} $60\%$;\ \pa{c} $180$.}}
\gap
\begin{probbox}{2}{\relax}
Một cửa hàng bán một chiếc xe đạp với giá $1\,800\,000$ đồng và được lãi $20\%$ so với giá vốn. Tìm giá vốn và số tiền lãi.
\end{probbox}
\probsol{
Giá bán bằng $100\%+20\%=120\%$ giá vốn.\par
Giá vốn: $1\,800\,000:120\times100=1\,500\,000$ đồng.\par
Tiền lãi: $1\,800\,000-1\,500\,000=300\,000$ đồng.
\dapso{Giá vốn $1\,500\,000$ đồng; tiền lãi $300\,000$ đồng.}}
\gap
\begin{probbox}{3}{\relax}
Một chiếc áo khoác giá $500\,000$ đồng. Cửa hàng giảm giá $10\%$, sau đó lại tăng giá mới thêm $10\%$. Hỏi giá cuối cùng là bao nhiêu? So với giá ban đầu, giá đã thay đổi bao nhiêu phần trăm?
\end{probbox}
\probsol{
Sau khi giảm: $500\,000\times90\%=450\,000$ đồng.\par
Sau khi tăng: $450\,000\times110\%=495\,000$ đồng.\par
Giá giảm $5\,000$ đồng, mà $5\,000:500\,000=1\%$.
\dapso{$495\,000$ đồng; giảm $1\%$.}}
\gap
\begin{probbox}{4}{\relax}
Một hình hộp chữ nhật có chiều dài $8$ dm, chiều rộng $5$ dm, chiều cao $4$ dm. Tính diện tích xung quanh và diện tích toàn phần của hình hộp đó.
\end{probbox}
\probsol{
Diện tích xung quanh: $(8+5)\times2\times4=104$ dm$^2$.\par
Diện tích một mặt đáy: $8\times5=40$ dm$^2$.\par
Diện tích toàn phần: $104+40\times2=184$ dm$^2$.
\dapso{$104$ dm$^2$; $184$ dm$^2$.}}
\gap
\begin{probbox}{5}{\relax}
Diện tích toàn phần của một hình lập phương là $294$ cm$^2$. Tìm cạnh và thể tích của hình lập phương đó.
\end{probbox}
\probsol{
Diện tích một mặt: $294:6=49$ cm$^2$.\par
Vì $7\times7=49$ nên cạnh dài $7$ cm.\par
Thể tích: $7\times7\times7=343$ cm$^3$.
\dapso{Cạnh $7$ cm; thể tích $343$ cm$^3$.}}
\gap
\begin{probbox}{6}{\relax}
Một bể nước dạng hình hộp chữ nhật dài $1{,}5$ m, rộng $0{,}8$ m, cao $1$ m. Mực nước trong bể bằng $75\%$ chiều cao của bể. Hỏi trong bể có bao nhiêu lít nước?
\end{probbox}
\probsol{
Chiều cao mực nước: $1\times75\%=0{,}75$ m.\par
Thể tích nước: $1{,}5\times0{,}8\times0{,}75=0{,}9$ m$^3=900$ dm$^3$.\par
$1$ dm$^3=1$ lít, nên trong bể có $900$ lít nước.
\dapso{$900$ lít.}}
\gap
\begin{probbox}{7}{\relax}
Một bể kính có đáy dài $50$ cm, rộng $40$ cm. Thả một hòn đá vào bể (hòn đá chìm hẳn trong nước) thì mực nước dâng lên $3$ cm. Tính thể tích hòn đá.
\end{probbox}
\probsol{
Thể tích hòn đá bằng thể tích phần nước dâng lên:\par
$50\times40\times3=6\,000$ cm$^3=6$ dm$^3$.
\dapso{$6\,000$ cm$^3$ ($6$ dm$^3$).}}
\gap
\begin{probbox}{8}{\relax}
Một khối gỗ hình lập phương cạnh $4$ cm được sơn tất cả các mặt, rồi cắt thành các hình lập phương nhỏ cạnh $1$ cm. Có bao nhiêu hình lập phương nhỏ được sơn \pa{a} $3$ mặt; \pa{b} $2$ mặt; \pa{c} $1$ mặt; \pa{d} không được sơn mặt nào?
\end{probbox}
\probsol{
Có tất cả $4\times4\times4=64$ hình lập phương nhỏ.\par
\pa{a} $8$ hình ở $8$ đỉnh được sơn $3$ mặt.\par
\pa{b} Mỗi cạnh trong $12$ cạnh có $4-2=2$ hình ở giữa: $12\times2=24$.\par
\pa{c} Mỗi mặt trong $6$ mặt có $2\times2=4$ hình ở phía trong: $6\times4=24$.\par
\pa{d} Phần lõi bên trong có $2\times2\times2=8$ hình. \ (\emph{Thử lại:} $8+24+24+8=64$.)
\dapso{$8$; $24$; $24$; $8$.}}
\gap
\begin{probbox}{9}{\relax}
Một ô tô khởi hành lúc $7$ giờ $15$ phút, đi với vận tốc $45$ km/giờ trong $2$ giờ $20$ phút. \pa{a} Ô tô đi được bao nhiêu ki-lô-mét? \pa{b} Ô tô đến nơi lúc mấy giờ?
\end{probbox}
\probsol{
\pa{a} $2$ giờ $20$ phút $=2\dfrac13$ giờ $=\dfrac73$ giờ.\par Quãng đường: $45\times\dfrac73=105$ km.\par
\pa{b} $7$ giờ $15$ phút $+\ 2$ giờ $20$ phút $=9$ giờ $35$ phút.
\dapso{\pa{a} $105$ km;\ \pa{b} lúc $9$ giờ $35$ phút.}}
\gap
\begin{probbox}{10}{\relax}
Hai thị xã A và B cách nhau $150$ km. Lúc $7$ giờ, một ô tô đi từ A đến B với vận tốc $40$ km/giờ, cùng lúc đó một xe máy đi từ B đến A với vận tốc $35$ km/giờ. Hỏi hai xe gặp nhau lúc mấy giờ và chỗ gặp nhau cách A bao nhiêu ki-lô-mét?
\end{probbox}
\probsol{
Sau mỗi giờ hai xe gần nhau thêm $40+35=75$ km.\par
Thời gian để gặp nhau: $150:75=2$ giờ, nên hai xe gặp nhau lúc $7+2=9$ giờ.\par
Chỗ gặp cách A: $40\times2=80$ km.
\dapso{Lúc $9$ giờ, cách A $80$ km.}}
\gap
\begin{probbox}{11}{\relax}
Lúc $6$ giờ, một người đi xe đạp từ nhà với vận tốc $12$ km/giờ. Lúc $7$ giờ, một xe máy đi từ cùng chỗ đó, trên cùng con đường, với vận tốc $36$ km/giờ. Hỏi xe máy đuổi kịp xe đạp lúc mấy giờ và chỗ đuổi kịp cách nhà bao xa?
\end{probbox}
\probsol{
Đến $7$ giờ, xe đạp đã đi trước $12\times1=12$ km.\par
Mỗi giờ xe máy gần xe đạp thêm $36-12=24$ km.\par
Thời gian đuổi kịp: $12:24=0{,}5$ giờ $=30$ phút, tức là lúc $7$ giờ $30$ phút.\par
Chỗ đuổi kịp cách nhà: $36\times0{,}5=18$ km.
\dapso{Lúc $7$ giờ $30$ phút, cách nhà $18$ km.}}
\gap
\begin{probbox}{12}{\stars{1}}
Hưng đạp xe lên dốc với vận tốc $12$ km/giờ và xuống dốc theo đường cũ với vận tốc $24$ km/giờ. Tính vận tốc trung bình của Hưng trên cả đoạn đường đi và về.
\end{probbox}
\probsol{
Coi con dốc dài $24$ km (độ dài bất kì đều cho cùng kết quả).\par
Lên dốc: $24:12=2$ giờ; xuống dốc: $24:24=1$ giờ.\par
Vận tốc trung bình: $(24+24):(2+1)=16$ km/giờ.
\dapso{$16$ km/giờ.}}
\gap
\begin{probbox}{13}{\relax}
Với vận tốc $40$ km/giờ, một xe khách đi hết quãng đường trong $3$ giờ. Muốn đi hết quãng đường đó trong $2{,}5$ giờ thì xe phải đi với vận tốc bao nhiêu?
\end{probbox}
\probsol{
Quãng đường: $40\times3=120$ km.\par
Vận tốc cần đi: $120:2{,}5=48$ km/giờ.
\dapso{$48$ km/giờ.}}
\gap
\begin{probbox}{14}{\relax}
Một chiếc thuyền đi xuôi dòng với vận tốc $18$ km/giờ, đi ngược dòng với vận tốc $12$ km/giờ. Tính vận tốc dòng nước và vận tốc của thuyền khi nước lặng.
\end{probbox}
\probsol{
Vận tốc xuôi dòng $-$ vận tốc ngược dòng $=2\times$ vận tốc dòng nước.\par
Vận tốc dòng nước: $(18-12):2=3$ km/giờ.\par
Vận tốc khi nước lặng: $18-3=15$ km/giờ.
\dapso{Dòng nước $3$ km/giờ; nước lặng $15$ km/giờ.}}
\gap
\begin{probbox}{15}{\stars{1}}
Một đoàn tàu dài $180$ m chạy qua một cột tín hiệu hết $12$ giây. \pa{a} Tính vận tốc của tàu theo m/giây và km/giờ. \pa{b} Tàu chạy qua một cây cầu dài $420$ m hết bao nhiêu giây?
\end{probbox}
\probsol{
\pa{a} Khi qua cột, tàu đi được quãng đường bằng chiều dài của nó: $180:12=15$ m/giây.\par
$15$ m/giây $=15\times3\,600:1\,000=54$ km/giờ.\par
\pa{b} Để qua cầu, tàu đi $180+420=600$ m: $600:15=40$ giây.
\dapso{\pa{a} $15$ m/giây $=54$ km/giờ;\ \pa{b} $40$ giây.}}
\gap
\begin{probbox}{16}{\stars{1}}
Một đường chạy hình tròn dài $400$ m. Hai bạn cùng xuất phát tại một điểm với vận tốc $3$ m/giây và $5$ m/giây. \pa{a} Nếu chạy ngược chiều nhau thì sau bao nhiêu giây hai bạn gặp nhau lần đầu? \pa{b} Nếu chạy cùng chiều thì sau bao nhiêu giây bạn chạy nhanh đuổi kịp bạn chạy chậm lần đầu?
\end{probbox}
\probsol{
\pa{a} Hai bạn cùng chạy hết một vòng: $400:(3+5)=50$ giây.\par
\pa{b} Bạn chạy nhanh phải hơn đúng một vòng: $400:(5-3)=200$ giây.
\dapso{\pa{a} $50$ giây;\ \pa{b} $200$ giây.}}
\gap
\begin{probbox}{17}{\relax}
\pa{a} Trong $5$ vòng thi, Ngân được $8$; $7$; $9$; $10$ điểm và một điểm nữa. Điểm trung bình của Ngân là $8{,}4$. Tìm điểm còn thiếu.\par
\pa{b} Tú tung một đồng xu $40$ lần, có $18$ lần xuất hiện mặt sấp. Viết tỉ số của số lần xuất hiện mặt sấp và tổng số lần tung dưới dạng phân số và dạng tỉ số phần trăm.
\end{probbox}
\probsol{
\pa{a} Tổng điểm $5$ vòng: $8{,}4\times5=42$.\par Điểm còn thiếu: $42-(8+7+9+10)=42-34=8$.\par
\pa{b} $\dfrac{18}{40}=\dfrac9{20}=45\%$.
\dapso{\pa{a} $8$;\ \pa{b} $\dfrac9{20}=45\%$.}}
\gap
\begin{probbox}{18}{\relax}
Lớp 5D có $40$ học sinh, trong đó $25$ bạn thích bóng đá, $18$ bạn thích cờ vua và $7$ bạn thích cả hai môn. Hỏi có bao nhiêu bạn không thích môn nào trong hai môn đó?
\end{probbox}
\probsol{
Số bạn thích ít nhất một môn: $25+18-7=36$ (vì $7$ bạn đã được đếm hai lần).\par
Số bạn không thích môn nào: $40-36=4$.
\dapso{$4$ bạn.}}
\gap
\begin{probbox}{19}{\stars{1}}
Mai dùng một nửa số tiền để mua sách, rồi dùng $\dfrac13$ số tiền còn lại để mua bút. Mai còn lại $40\,000$ đồng. Hỏi lúc đầu Mai có bao nhiêu tiền?
\end{probbox}
\probsol{
Tính ngược: $40\,000$ đồng ứng với $1-\dfrac13=\dfrac23$ số tiền còn lại sau khi mua sách.\par
Sau khi mua sách còn: $40\,000:\dfrac23=60\,000$ đồng.\par
Đó là một nửa số tiền, nên lúc đầu Mai có: $60\,000\times2=120\,000$ đồng.
\dapso{$120\,000$ đồng.}}
\gap
\begin{probbox}{20}{\stars{1}}
Một trang trại có gà và thỏ, đếm được $35$ cái đầu và $94$ cái chân. Hỏi mỗi loại có bao nhiêu con?
\end{probbox}
\probsol{
Giả sử tất cả là gà: $35\times2=70$ chân, thiếu $94-70=24$ chân.\par
Mỗi con thỏ nhiều hơn $2$ chân: số thỏ $24:2=12$ con; số gà $35-12=23$ con.
\dapso{$23$ con gà, $12$ con thỏ.}}
